\chapter{}
Muitas pessoas indagam, nos dias de hoje, quando já não é comum encontrar famílias com cinco filhos, se foi difícil dar conta de criá-los.
Eles próprios, hoje pais e mãe, perguntam-me isso, vez ou outra.
Afirmo que não foi, porque não me ficaram lembranças de momentos especialmente difíceis, exceto talvez aqueles em que tive que me preocupar com a saúde deles.


Afora as ocasiões de doença, sempre digo que me parecia, às vezes, que os pais que eu conhecia de um ou dois filhos pareciam ter muito mais trabalho que nós.
É fato que quando eram muito pequenos, fui ajudada pela proximidade dos meus pais, minhas tias, meus irmãos e cunhados e pela existência do bendito sítio.
Paulo, por outro lado, apesar de ter que se ausentar por longos períodos, sacrificava-se para manter duas empregadas em casa e, quando vinha, a cada mês, procurava aliviar-me de toda a maneira, além de mostrar-se extremamente compreensivo com o tempo que as aulas, as correções, a doença do meu pai, e as cinco crianças me tomavam.
Houve uma época em púnhamos nossas conversas em dia a partir das dez horas da noite, quando íamos para a cozinha juntos, preparar as vinte e seis mamadeiras que ficavam armazenadas na porta da geladeira para o dia seguinte.


Porém, a ajuda mais importante mesmo, veio das próprias crianças.
Tenho para mim, e sempre propunha isso aos pais em palestras, que filhos devem ser educados como sócios nesse seriíssimo empreendimento que é a edificação de uma família.
Claro que, no meu caso, isso não foi pensado a priori, mas surgiu da própria necessidade.
Todos ajudavam a tomar conta de todos.
E dos cachorros, quando o pai cismou que devíamos ter um monte deles.
Eles próprios intuíam que, sendo eu uma só para atender a tudo, não havia espaço para muito personalismo, muito individualismo, muita manha.
E que disciplina e organização eram importantes para a sobrevivência do grupo.
Então, eu nunca soube o que era me afligir com menino que não comia ou não dormia às horas certas.
Muito comumente eu tinha que levar os cinco comigo ao supermercado e não havia choradeira para que eu comprasse isso ou aquilo.
Na maior parte das vezes, aceitavam de boa vontade o revezamento à janela do carro e por muitos anos viajaram milhares de quilômetros, por horas a fio, cruzando esse país para lá e para cá, os cinco, sem dar maiores alterações.
É surpreendente lembrar a tolerância que demonstravam nessas ocasiões.


O que gosto de pensar da experiência de educá-los, é que foram crianças à moda de antigamente, frutos amadurecidos no pé.
Se houve algum mérito da minha parte e do Paulo, além da segurança afetiva que não lhes faltou, foi o de nos esforçarmos para que vivessem cada etapa da vida no tempo e na hora certos, com calma e intensidade, responsabilizando-os sempre por aquilo que estava à altura deles resolver, mas tentando a todo custo evitar-lhes situações que ainda não estivessem preparados para enfrentar.


\begin{figure}[H]
\centering
\includegraphics[width=\linewidth]{36/tomaz-rescan.jpg}
\caption*{Tomaz.}
\end{figure}

Tomaz, competente, reservado e disciplinado, era apontado como modelo pelos professores e funcionários das escolas por onde passou.
Na verdade, nele se somavam caráter e personalidade à condição de mais velho dos cinco.
Assumiu, desde muito cedo, a responsabilidade de dar exemplo e ajuda na educação dos irmãos.
Por outro lado, a julgar pela minha própria experiência como primogênita, o sacrifício que essa situação quase sempre implica parece provocar em resposta um intenso desejo de autonomia.
No primeiro dia em que conduzi Tomaz ao jardim da infância, assim que soou o sinal de entrada, sério e compenetrado, ele caminhou em direção à sala de aula, carregando sua mochilinha nos ombros, sem nem olhar para trás, deixando-me ali um tanto desconcertada em meio a um mar de crianças recalcitrantes que se esgoelavam agarradas às saias das mães.
Uma dúzia de anos depois, quando Paulo e eu o levamos ao aeroporto de Teresina para embarcar rumo a São Paulo, onde prestaria o vestibular, repetiu-se a cena: puxando a mala atrás de si, sério e compenetrado, vimo-lo atravessar a pista e embarcar no avião sem olhar para trás, ainda que todos os três soubéssemos que, a partir dali, ele começava a trilhar o seu próprio caminho.
Consolou-me Padre Florêncio, sempre ele:

\textit{``-- Filha, nada receie por esse menino porque o que quer que ele escolha fazer nessa vida, esteja certa de que fará bem feito.''}

Hoje, há quem ache que Tomaz parece um sujeito muito centrado em si mesmo, o que ele é, sem dúvida, meio ermitão e sistemático, o que também ele é, mas de uma coisa tenho certeza: é das raras pessoas que promove e respeita a autonomia do outro e, quando ajuda, o faz na medida certa, na hora justa e com extremo cuidado pela sensibilidade daquele que precisa.
Curiosamente, essas eram as qualidades que eu mais admirava em meu pai.


\begin{figure}[H]
\centering
\includegraphics[width=\linewidth]{36/marcelo-rescan.jpg}
\caption*{Marcelo.}
\end{figure}

Marcelo, como cabe a um segundo filho, sempre foi extrovertido, imaginativo e dono de um insopitável senso crítico, vazado sob um humor rápido e certeiro e um sorriso irresistível.
Por isso mesmo criou-me, em duas ocasiões, problemas com professores mais susceptíveis que, felizmente para mim, acabaram rendidos diante da sua reconhecida facilidade para alcançar notas máximas.
Aliás, ele sempre deu às pessoas a impressão de que para ele, tudo acontecia com facilidade.
Todos os cinco lograram passar nos vestibulares para as melhores faculdades sem frequentar os famigerados cursinhos, mas Marcelo, que escolhera Engenharia Mecânica, saiu-se como sempre tão bem que acabou sendo convidado a fazer Engenharia Mecatrônica, modalidade recém-criada na Universidade de São Paulo, e que nascera com o propósito de reunir os melhores alunos da casa.
Tenho certeza de que ele não gastou nem dois terços do tempo de estudo e muito menos do custo em mensalidades despendidos pela privilegiada elite dos distinguidos pelo convite.
Tanto é que, para minha secreta satisfação, ele descobriu-se bem pouco afinado com a maioria dos robozinhos competitivos e preconceituosos da sua sala e acabou indo buscar seus amigos em ambientes mais humanizados da escola.
Na verdade, sua empatia não só tornou-o um bom fazedor de amigos como fez dele, depois, um líder talentoso no trabalho.
Ao lado de tudo isso, Marcelo sempre exibiu certo pendor dramático.

Tanto ele como Tomaz, durante os primeiros anos de faculdade, tiveram que enfrentar tempos heroicos, porque em casa atravessávamos um período de dificuldades sérias.
Com Marcelo, entretanto, o aperto se revestiu de cores barrocas, para a diversão do tio João que o acudiu no transe, Mas, justamente essa qualidade de expor o que sente com rara sinceridade é que o levou a proporcionar, àquela que hoje é sua mulher e ao resto da família, um momento de emoção inesquecível - a bela elegia que dedicou à noiva na festa do casamento de ambos foi de tal modo tocante, que pôs em lágrimas até os machos mais empedernidos ali presentes.
E é ainda essa natureza, que nele faz conviverem em harmonia inteligência e emoção, que faz dele, hoje em dia, um pai carinhoso e devotado ao extremo.


\begin{figure}[H]
\centering
\includegraphics[width=\linewidth]{36/fernando-rescan.jpg}
\caption*{Fernando.}
\end{figure}

Fernando demorou a crescer e soube tirar vantagem disso no colégio.
Observador, quietinho e míope, despertava o instinto maternal das professoras e das colegas.
Mas a ala masculina não se deixava enganar.
Debaixo daquela fofura se escondia uma lógica, um poder de articulação que fazia dele o ideólogo predileto da turma na organização de brincadeiras e eventos e de outras iniciativas nem tão inocentes assim, mas que raras vezes chegavam aos meus ouvidos na Coordenação.
Aprendia com facilidade e rapidez o que lhe deixava tempo de sobra para dar vazão à inventividade.
Um professor que eu reputava o melhor conhecedor dos nossos alunos me disse mais de uma vez ter certeza de que Fernando estava no centro de todas as ``aprontações'' da turma, mas que surpreendê-lo em ação era dificílimo.
Por mais rápido que fosse na tentativa, tudo que conseguia flagrar era aquele olhar de absoluta inocência na carinha séria.

Alguém pode pensar que, nascido sob o signo de Coelho no horóscopo chinês, Fernando seja um dissimulado.
Não é.
Apenas gosta de surpreender-nos com o resultado daquilo que burilou cuidadosamente em sossego, como uma obra de arte.
Que pode ser um evento, um texto contundente, uma pintura, um prato elaborado.
Ou uma palpitante aventura como aquela, secretíssima, perpetrada em combinação com seu amigo Rachib, nos tempos do colégio, em que uma excursão de bicicleta pela mata dos arredores de Teresina acabou num formidável tombo que o atirou para cima de um tronco de tucum.
Espetado de alto a baixo como um ouriço e coerentemente caladinho, ele aguentou por uma tarde inteira eu arrancar da sua carne dolorida, com pinça de sobrancelha, os insidiosos espinhos, um a um.
Ou ainda, uma demonstração sobre patins, feita para uma plateia de alunos, pais e professores reunidos numa quadra de ginásio para a movimentada gincana anual da escola, durante a qual, após uma série de quedas alucinantes que me fizeram perder a compostura de coordenadora e berrar desesperadamente para que o detivessem antes que se matasse, ele saiu deslizando como um astro do esporte, sob aplausos retumbantes, deixando-me com a desconfortável impressão de que minha histeria contribuíra para dar brilho ao espetáculo.
Eu nunca soube que ele tivesse aprendido a patinar.

Com o tempo, também nos daríamos conta de que Fernando puxou para si boa parte do que havia de apreço às boas coisas da vida na herança genética da família.
Embora, para minha satisfação, todos os cinco saibam valorizar o que há de bom e bonito no mundo, Fernando sempre demonstrou uma firme persistência em fazer disso um propósito de vida.
Quando, depois de formado na faculdade, ele foi morar na Europa, descobri, entre as coisas que deixou para trás, um bauzinho de vime onde, desde menino, ele armazenou cartões postais, prospectos e recortes de paisagens do mundo inteiro solicitados através de cartas enviadas a consulados e embaixadas, ao lado de belos \textit{folders} de divulgação de queijos, vinhos e embutidos finos da França e da Itália.
Era, com certeza, uma espécie de culto propiciatório ao que ele esperava alcançar no seu futuro.
E que aos poucos, do seu jeito, silenciosa e caprichosamente ele vai construindo.
Como uma obra de arte.


\begin{figure}[H]
\centering
\includegraphics[width=\linewidth]{36/paula-rescan.jpg}
\caption*{Paula.}
\end{figure}

Paula foi o mais risonho bebê de casa.
Ria até regurgitar, o que me fez andar por uns bons meses recendendo a queijo de manhã até a noite.
Mas, acho que logo se deu conta de que precisaria endurecer para se impor numa família onde ela e eu estávamos em franca minoria.
E assim, pouco a pouco, começou a exibir uma surpreendente intransigência ao defender suas posições, ao que ela acrescentou, com o tempo, um arzinho professoral de quem sabe tudo.


Quando ela nasceu, foi uma enorme surpresa.
Por três anos seguidos, eu chegava à maternidade munida de coloridos laços de fita à espera da menina que tanto o pai queria e saía de lá com mais um machinho nos braços.
Quando, na quarta tentativa, ela apareceu, senti-me insegura.
Tinha medo de não saber mais criar uma filha.
Talvez por isso, empenhei-me desde cedo em realçar sua feminilidade.
Enchi-a de bonecas, panelinhas e, um pouco mais tarde, por muitos meses ela pode ser vista carregando de um lado para o outro, sem muito entusiasmo, uma cestinha de bordadeira, linhas coloridas e uma tela sobre a qual deveria surgir um leão ou um pavão, não me lembro, que jamais se materializou.
Era minha tentativa de despertá-la para as artes manuais, nas quais, aliás, sempre fui um completo fracasso.
Mas Paula, para desespero dos irmãos, preferia engajar-se nas brincadeiras deles e, verdade seja dita, eles tinham com ela uma paciência infinita, embora vez ou outra se vingassem, rolando de rir quando ela assomava à porta do banheiro vestida em robezinho e chinelos combinando, carregando a ``necessaire'' e trazendo os cabelos escondidos sob uma touca de banho de babadinhos.

No ``Sossego da Mamãe'', escolinha onde ela e Vicente estudaram em Conceição do Araguaia, Paula, embora não se furtasse às atividades, parecia um tanto arredia.
Mas, ao chegar ao Diocesano, apesar do estranhamento inicial, pareceu desabrochar, enturmou-se sem dificuldade e fez-se querida dos professores que a consideravam exemplo de responsabilidade, alegria e gentileza.
Até mesmo o temível Padre Florêncio, que preferia ver as meninas aprendendo a cozinhar e cuidar da família derretia-se por ela.
Fez muitas amigas e amigos entre os colegas, virou festeira e tanto entrosamento acabou por me criar dificuldades.
Tive que me opor a certos usos correntes entre as moças do lugar.
Entre eles, o das meninas exibirem-se precocemente em decotes e saias curtas e justas, prática a que o calor servia de justificativa.
É verdade que quando viajávamos para São Paulo eu constatava que aqui no sul os hábitos não pareciam melhores, porém há nas mulheres nordestinas uma faceirice natural, uma desinibição e um gosto por enfeites e atavios que lhes confere uma feminilidade ostensiva e provocante, ausente nas paulistas, no geral mais apagadinhas.
Convencer Paula a vestir-se de modo condizente com a sua pouca idade foi luta ferrenha que atraiu em alguns momentos até a intervenção de Paulo, a favor da filha, claro.
Porque minha única garota sabia ser teimosa na defesa dos seus pontos de vista até o limite do insuportável.
A gritaria foi grande.
Hoje, ela ri quando se lembra de como se sentia nas festinhas, vestida de menina-moça comportada.
Mais ou menos como uma Rainha Elizabeth dançando entre ``tchutchucas'' num baile funk.
Aliás, coitada, como \textit{boutiques} e grifes não cabiam em nosso orçamento, ao chegarmos para as férias em Araraquara, ela de novo tinha que enfrentar discriminação, agora das primas, cujo bem-estar não dispensava tais refinamentos.
Mas, diga-se a favor da nossa filha que ela parecia nem se dar conta dos olhares preconceituosos e, se dava, parecia saber atribuir-lhes o enquadramento correto.

Se eu algum dia temi não ser capaz de ajudar uma filha a construir-se como um ser humano de qualidade, ela soube dissipar todos os meus medos.
Paula amadureceu e se fez mulher no seu tempo e no seu ritmo, lento e tranquilo.
Hoje eu e ela temos uma convivência de velhas amigas, daquelas capazes de rir muito juntas.
Isso para mim tem um encanto extraordinário.
Paulo e eu costumamos comentar, com orgulho, que a fisionomia aberta e franca dela sempre a destaca, onde quer que ela esteja e lhe confere uma beleza muito pessoal que independe de etiquetas e cosméticos.
E, sempre do seu jeito, vai passando pelos desafios da vida com tal fleuma que às vezes é preciso parar para observar quanto essa menina tem sido capaz de fazer.
Como, por exemplo, viver a maternidade com a calma e a segurança de uma veterana.
Minha filha, percebo com satisfação, está pronta para ir muito mais longe do que fui capaz de chegar.
Só não lhe cobrem pressa.

Todo pai deve agradecer aos céus quando os filhos mais velhos se encaminham bem na vida, pois, via de regra, é certo que os mais novos lhes seguirão as pegadas.
Paula e Vicente, como os irmãos, não deram o menor trabalho na escola, ou fora dela.
Mas, não teve como evitar, arcaram, eu suspeito, com a obrigação auto-imposta de não fazer menos do que os mais velhos fizeram.
Isso nunca lhes foi cobrado, por isso me surpreendi quando Paula, sabe-se lá de onde, sacou a ideia de prestar Medicina na Universidade de São Paulo, o vestibular mais concorrido do país.
Tudo bem, exceto pela evidência de que nada nela combinava com tal escolha.
Aparentemente, era só o receio de parecer menos capaz se optasse por uma carreira menos ``difícil''.
Felizmente ela não alcançou a classificação necessária, mas, de frustração, engordou uns bons quilos que lhe custaram muito esforço perder.
O episódio me ajudou a entender melhor as reações do Vicente quando a hora dele chegou.

\begin{figure}[H]
\centering
\includegraphics[width=\linewidth]{36/vicente-melhor.jpeg}
\caption*{Vicente.}
\end{figure}

Vicente já nasceu pensando o mundo e pela carinha com que olhava para tudo e para todos, parecia achar que muita coisa podia ser melhorada.
Observando-o, minha tia Maria Angelina dizia:

\textit{``-- Esse menino vai ser filósofo, com certeza.''}

\begin{figure}
\centering
\includegraphics[width=\linewidth]{36/vicente-escoteiro-melhor.jpg}
\caption*{Vicente, escoteiro.}
\end{figure}

Cresceu saudável, foi o único que não passou pelas habituais perebas da infância e deve ser por isso que até hoje tem pavor de ficar doente e de tomar injeção.
Gregário, sempre gostou de ver-se cercado de gente, de modo que não achou graça alguma em ver-se arrancado do convívio dos tios e primos para ser levado, ainda em fraldas, para Conceição do Araguaia.
Embora jamais tenha sido birrento ou manhoso, passou um bocado de tempo emburrado.

Seguia os irmãos por toda parte e não admitia ser deixado para trás só por que era pequeno.
Foi assim que, aos poucos, foi aceitando a ideia de ficar longe de Araraquara em troca de andar de canoa, ir com eles e o pai para a fazenda, encarapitar-se no trator, pescar e banhar-se no rio Araguaia.
Só uma vez chorou para valer: foi quando, no desfile da escolinha maternal, pelo 7 de setembro, não lhe deram a fantasia de tocador de tambor e o incluíram no bloco dos palhacinhos, o que ele claramente tomou como ofensa à sua dignidade.

Como os demais, foi muito bom aluno.
Só não passou por Padre Florêncio porque quando ia para o primeiro ano do curso colegial, nos mudamos de volta para Araraquara.
Eu já lamentava antecipadamente ter que recorrer a um cursinho para cobrir as conhecidas deficiências das escolas da minha cidade, mas Vicente não perdeu qualidade.
Ainda que fizesse tudo para passar despercebido, era fácil identificar nele os atributos que o tornariam um competente engenheiro no futuro.
Apesar disso, escaldada com a experiência que tivera com a Paula, sugeri a ele confirmar essa inclinação através de uma sondagem vocacional.
O resultado não surpreendeu ninguém.
Surpreendida ficou a psicóloga, segundo ela própria confessou, ao se deparar com alguém, da geração dele, capaz de priorizar valores como família, amizade e qualidade de vida, acima de riqueza, poder ou prestígio.

Difícil com Vicente foi conseguir que ele saísse da casca.
O mais araraquarense dos meus filhos sempre detestou se colocar em evidência.
Atrás da modéstia se escondia, na verdade, uma autocrítica feroz.
Meu Torquemada, como eu o chamava, era muito severo consigo próprio.
Nunca se achava bom o suficiente.
Dava a impressão de que nunca aceitaria nada menos que a perfeição como meta para si, para os outros e para o mundo.
Isso rendeu horas de discussão em que eu tentava convencê-lo de que a Vida jamais se submete aos rigores da prancheta.
Só hoje, homem feito, parece começar a aceitar essa evidência.
Só queria que o próximo passo fosse descobrir que nisso consiste toda a graça de viver.
